% Structural refactor 2026-08-27. Historical source: Team 7.
\chapter{Gold Across Monetary Regimes and the Debasement Narrative}
\chapterstatus{Historical material is edited reuse from Team 7. The post-COVID and
debasement discussion is an institutional-source bridge. The debasement trade is
treated as a contested market narrative, not as an established causal law.}
\sourceassetnote
\section{Gold as a Monetary Anchor}
\subsection{The Classical Gold Standard}
What is commonly referred to as the classical gold standard --- a commitment by participating countries to fix the prices of their domestic currencies in terms of a specified amount of gold, with national money and other forms of money (bank deposits and notes) freely convertible into gold at the fixed price --- governed the western economy roughly between 1870 and 1914. It was not the outcome of a single treaty or international conference. Rather, it emerged gradually from the historically constrained choices of national governments. To understand its genesis, it is necessary to look at the bimetallic system that preceded it, in which both gold and silver circulated as legal tender.
 
Silver had long been the dominant money in medieval and early modern Europe, because it could be minted into coins of a value convenient for everyday transactions. Gold, by contrast, was too scarce and too valuable, making gold coins unsuitable for routine commerce. Improvements in minting technology, such as the spread of steam-powered machinery, were a necessary condition for the technical feasibility of a pure gold standard, but even then gold coins had to be supplemented by token coins and, in the end, by paper notes whose value rested on the promise of conversion into gold on demand.\sourcecitation{Eichengreen1996}
 
To maintain the bimetallic system, governments were required to fix a mint ratio between the two metals; for example, in France in 1803 a gram of gold was worth 15.5 grams of silver (the ``loi du 7 germinal, an XI''). In practice it was difficult to keep both metals in simultaneous circulation. As suggested by Gresham's law, which states that ``bad money drives out good'', whenever the market price of the metals deviated from the official mint ratio arbitrageurs would hoard or export the metal undervalued by the mint and supply the overvalued one for coinage, so that only the overvalued metal remained in circulation. The first transition to a de facto gold standard arose from this logic. In 1717, Sir Isaac Newton, who served as Master of the Mint for Great Britain, fixed the gold guinea at 21 silver shillings, with the effect of overvaluing gold relative to silver. Silver coins disappeared from circulation as they were melted down for bars, leaving gold as the principal circulating medium.\sourcecitation{BordoKydland1992}
 
After 1717, despite Britain's dominance, the emergence of a global gold standard was far from immediate. In the first half of the nineteenth century, France anchored a bimetallic bloc that stabilised global metal prices and acted as a bridge between silver-standard nations --- such as the Austro–Hungarian empire, the Scandinavian countries and the various German states --- and gold-standard Britain. Following the Franco–Prussian war (1870–1871), France had to pay a massive indemnity of five billion francs to Germany. The newly unified German empire used these resources to overhaul its monetary system, replacing the set of old silver currencies with the new gold-backed mark.
 
Germany's move flooded the global market with demonetised silver, driving down its price. This forced the countries in the Latin Monetary Union, a unified system of coinage founded in 1865 by France, Italy, Belgium and Switzerland, to suspend silver coinage in order to avoid inflation, effectively pushing them onto gold. More importantly, the birth of the mark unleashed strong network externalities. Great Britain and Germany were at the time the dominant powers in trade and finance, so their joint adoption of the gold standard made it overwhelmingly attractive for other nations to follow. Sharing a monetary standard with powerful commercial partners reduced exchange-rate uncertainty, lowered transaction costs and, above all, cut borrowing costs in the London capital market. A gold peg acted as a ``good housekeeping seal of approval'', signalling fiscal prudence to foreign investors.\sourcecitation{BordoRockoff1996} By 1900, the gold standard had encompassed most of the global economy, including Russia, Japan and the United States.\sourcecitation{Eichengreen1996}
\subsection{Convertibility and Fixed Parities}
The technical foundation of the gold standard was the commitment of each central bank in the participating countries to buy and sell its domestic currency for a fixed quantity of gold. As a consequence, exchange rates between currencies were effectively fixed by the ratio of the weight of gold in one currency to the weight of gold in another --- the so-called mint parity. Fluctuations in market exchange rates were limited within narrow bands, the gold points, determined by the costs of physically shipping gold between financial centres. This implied that the system operated differently across countries, depending on their location and status.
 
The eighteenth-century Scottish philosopher David Hume articulated the theoretical mechanism by which the system was supposed to maintain equilibrium. His price–specie–flow model described how a country running a balance-of-payments deficit would experience an outflow of gold, reducing its money supply and lowering its domestic price level. Falling prices would make its exports more competitive and its imports more expensive, gradually reversing the deficit and limiting the gold outflow; surplus countries would experience the opposite process. In theory, the system was therefore self-equilibrating.
 
In reality, as stressed by modern economic historians such as Eichengreen,\sourcecitation{Eichengreen1996} central banks played an active role in shaping domestic monetary conditions. Their main tool was interest-rate policy: raising rates to attract foreign capital when reserves fell and lowering them when gold was abundant. The true channel of adjustment was therefore as much the international flow of short-term capital as the physical movement of gold envisaged by Hume. In this context, the position occupied by the Bank of England was crucial. As the manager of the most important currency for international trade and the banker to the world's leading creditor nation, an adjustment in its discount rate had the power to influence global monetary conditions. For example, a rise in the Bank Rate --- the interest rate the central bank charges commercial banks for loans --- would draw gold and capital to London from across the world, stabilising sterling without the deflation implied by Hume's simple mechanism.
 
The convertibility commitment was also a political way for governments to signal their willingness to subordinate domestic economic concerns to the requirements of external balance. Market trust in this commitment tended to be self-reinforcing. When a currency came under pressure, investors expected adjustment from the central bank rather than devaluation; instead of selling the currency, they often bought it, thereby supporting their own expectations.
 
In core countries such as Britain, France, Germany and the United States, a pure gold standard operated with high credibility. In the ``periphery'' of the world economy (Latin America, Eastern Europe and parts of Asia), convertibility was much more fragile. These nations mainly exported raw commodities and were exposed to frequent terms-of-trade shocks. Their shallow domestic financial markets forced them to borrow in foreign currencies and, to economise on scarce physical gold, many adopted a ``gold-exchange standard''. They held foreign-exchange reserves in their central banks, rather than gold --- mainly sterling balances, considered ``as good as gold''.\sourcecitation{EichengreenEtAl2003}
\subsection{Gold and Monetary Discipline}
The classical gold standard offered a politically insulated discipline on the expansion of money and credit, something that no purely discretionary monetary regime could provide. The quantity of money was tied to gold reserves, which could grow only as fast as mines were able to extract new metal, making it impossible for governments to print their way out of fiscal difficulties. In this sense the gold standard bound the hands of politicians, acting as a constraint on fiscal profligacy.\sourcecitation{BordoKydland1992}
 
The system's functioning relied on specific conditions. First, it presupposed an intellectual climate in which the maintenance of the gold standard was regarded as a priority, overriding other claims such as the pursuit of full employment. Second, it required open and flexible markets, capable of absorbing the adjustments demanded by the system without causing drastic social consequences for labour, goods and capital. Third, it depended on a high level of international cooperation among central banks to manage crises before they became systemic. As mentioned above, London played a central role as a source of discipline and cohesion: the Bank of England's discount rate worked as something close to a world interest rate.\sourcecitation{Eichengreen1996}
 
By the end of the First World War, these conditions were already being eroded. The diffusion of unionism and the bureaucratisation of work made wages less responsive to downward pressure than gold-standard adjustments required, while organised labour was gaining a political voice that clashed with governments prioritising exchange-rate stability over employment. At the same time, the financial system became more efficient thanks to fractional-reserve banking but also more fragile, creating new vulnerabilities to confidence crises that gold-standard rules were unable to address. As a result, the institutional environment that had supported the classical gold standard did not survive the war.
\section{The End of Gold Parity}
\subsection{Bretton Woods and the Dollar}
During the interwar decades, the attempt to restore the classical gold standard proved inadequate. In 1925 Britain returned to gold at the pre-war parity, but its position had changed profoundly during the war years. It had been forced to liquidate much of its overseas investment to finance the conflict, permanently weakening London's role as the world's banker. At the same time, the United States was emerging as the leading creditor nation, yet the political and institutional conditions were not yet in place for it to assume the responsibilities that Britain's lost hegemony entailed.

The war had made international cooperation far more fragile. Prices and wages had become less flexible, and governments more sensitive to unemployment. As a consequence, the gold-exchange standard that emerged in the 1920s --- based on reserves of gold, sterling and dollars --- created new vulnerabilities. At the base of the system lay confidence in the two reserve currencies; when that confidence wavered, as it did after the Great Depression, the system rapidly fell apart. Between 1931 and 1936 almost every country abandoned gold to free its monetary policy and avoid severe deflation. Although the manner and timing of this choice varied, the main consequence was a return to competitive ``beggar-thy-neighbour'' devaluations and draconian trade quotas that throttled global commerce.

To design a new system --- one that preserved the stability of fixed exchange rates while granting countries enough flexibility to pursue domestic full employment --- 44 nations gathered at Bretton Woods, New Hampshire, in 1944. Two main visions were debated. Representing Britain, John Maynard Keynes proposed a global clearing union that would issue a new reserve asset, the ``bancor'', with generous overdraft facilities to help countries sustain full employment without resorting to deflation, and a symmetric obligation on both surplus and deficit countries to adjust. On the other side, the US representative Harry Dexter White envisaged an International Monetary Fund that would supervise a system of pegged exchange rates and provide short-term balance-of-payments support. White's design left the burden of adjustment primarily on deficit countries and preserved the central role of the dollar.\sourcecitation{Eichengreen1996,ObstfeldTaylor2004}

The outcome was largely White's design, with one important addition: the adjustable peg, under which countries could devalue their currencies to correct a ``fundamental disequilibrium''. Every country was required to declare a par value for its currency in terms of gold, or of a currency directly convertible into gold. By 1945 the United States held over 70\% of the world's monetary gold, so the system was a de facto gold–dollar standard. The US pegged the dollar at \$35 per ounce, and the rest of the world pegged to the dollar.\sourcecitation{Mishkin2019}
\subsection{Tensions in Convertibility}
Through the 1950s and 1960s the Bretton Woods system maintained a remarkable period of exchange-rate stability during years of unprecedented economic growth in Western Europe and Japan and of expanding international trade.Nevertheless, its architecture contained a fatal structural flaw, identified by the economist Robert Triffin in 1960 and since known as the Triffin dilemma. The growth of the world economy expanded the demand for international reserves, which under Bretton Woods relied on the United States running balance-of-payments deficits to supply the necessary dollar liquidity. But while the pool of dollars held by central banks had to grow in step with the world economy, the stockpile of physical gold in Fort Knox was finite. As a result, the United States' capacity to honour its commitment to convert dollars into gold at \$35 per ounce steadily deteriorated, making an eventual crisis of confidence inevitable.\sourcecitation{Triffin1960,Eichengreen1996}

During the 1960s, Triffin's prophecy began to materialise. In October 1960 the price of gold on the private market spiked above \$40 per ounce, well above the official \$35 --- a clear signal of market scepticism about dollar convertibility. The United States deployed a series of stopgap measures to defend the peg without deflating its economy. In 1961 the ``London Gold Pool'' was formed, an agreement among the US and European central banks to sell bullion aggressively into the private market in order to cap the price at \$35.

These operations ultimately proved insufficient. In June 1967 France withdrew from the Gold Pool, in protest at what it regarded as America's ``exorbitant privilege'' --- a phrase coined by Finance Minister Val{\'e}ry Giscard d'Estaing and embraced by President Charles de Gaulle, who had already spent years converting France's dollar reserves into gold and repatriating the bullion to Paris. By March 1968 the Gold Pool had collapsed, and convertibility for private actors came to an end.
\subsection{The Nixon Shock and Fiat Money}
The war in Vietnam and the Great Society programmes pushed the US balance of payments into an even worse position. By 1970 inflation had climbed toward 6\%, and in 1971 the United States recorded its first trade deficit of the twentieth century. Under the system of pegged exchange rates, this set off a mechanism that led European countries --- West Germany above all --- to ``import'' American inflation. In May 1971, judging the situation politically unacceptable, West Germany abandoned its dollar peg and let the mark float upward. The Swiss and the Dutch also floated, and speculative pressure on the dollar intensified.

On 15 August 1971, US President Richard Nixon announced on national television the suspension of the dollar's convertibility into gold, a 10\% surcharge on imports, and a 90-day freeze on wages and prices. The announcement --- which became known as the ``Nixon Shock'' --- came without prior consultation with America's allies and brought to an end the Bretton Woods system that had governed the world economy for the previous quarter-century.

An attempt to salvage the fixed-rate system was made by the major powers at the Smithsonian Institution in Washington in December 1971. The dollar was devalued to \$38 per ounce of gold, other currencies were revalued against it, and the margins within which exchange rates could fluctuate were widened. The new arrangement, which Nixon welcomed with great enthusiasm, lasted barely fourteen months. After further speculative pressure, in March 1973 the major currencies moved to floating exchange rates.

This transition --- to floating exchange rates and fully fiat money --- marked a decisive rupture in monetary history. For the first time in the modern international economy there was no metallic anchor, no convertibility commitment, and no automatic mechanism constraining the expansion of money and credit. Everything now depended on the wisdom and credibility of central banks and governments, on the quality of monetary institutions rather than on the discipline of gold.\sourcecitation{Eichengreen1996,Mishkin2019}
\section{Gold and Silver as Financial Assets}
\subsection{Safe Haven and Hedge Roles}

Gold and silver are often described as ``safe'' assets, but the precise meaning of this
label is more nuanced in the portfolio and risk–management literature. A standard
distinction is made between \emph{hedges}, \emph{safe havens}, and \emph{diversifiers},
each capturing a different form of protection against adverse market movements (see, e.g.,
\sourcecitation{BaurLucey2010,ErbHarvey2013}). Clarifying these concepts is essential before
turning to the empirical analysis of co–movements with bonds, equities, and the US dollar.

A \emph{hedge} is an asset that is negatively correlated, or at least uncorrelated, with
another asset or portfolio on average, i.e., over normal market conditions. In this sense,
an asset qualifies as a hedge for equities if its returns tend to move in the opposite
direction of stock–market returns, or if it shows no systematic co–movement with them in
typical periods. A \emph{strong hedge} would exhibit a persistently negative correlation,
whereas a \emph{weak hedge} might be characterized by correlations that are close to zero
but not significantly positive \sourcecitation{BaurLucey2010}.

A \emph{safe haven}, by contrast, is defined with respect to extreme market conditions
rather than average behavior. An asset is said to act as a safe haven for equities if it
is uncorrelated or negatively correlated with stock returns specifically during periods of
market stress, such as crashes, crises, or severe drawdowns. In tranquil times the
asset may even be positively correlated with risky markets; what matters for the safe–
haven property is the ability to preserve or increase value when investors most need
protection. This notion naturally calls for a conditional state–dependent analysis of
co–movements, focusing on the tails of the return distribution and on crisis subperiods
\sourcecitation{BaurLucey2010}.

A third category is that of \emph{diversifiers}. A diversifier is an asset that is
positively but not perfectly correlated with a given portfolio in normal times. Although
it does not hedge losses one–for–one, it still improves the risk–return trade–off by
lowering overall portfolio variance through imperfect co–movement. Many assets fall into
this category: they are not true hedges or safe havens, but they contribute to portfolio
stability when combined with other holdings.

Within this classification, gold is often argued to behave as a hedge or even as a safe
haven against equity and currency risk, particularly during episodes of heightened
uncertainty and financial turmoil \sourcecitation{BaurLucey2010,ErbHarvey2013}. Silver has sometimes
been considered a weaker safe asset, with stronger ties to industrial demand and
business–cycle conditions. Whether these characterizations hold empirically, and under
which market regimes, is an open question that motivates the correlation and regression
analysis carried out in later chapters.

In the framework of this paper, we interpret hedge and safe–haven properties through the
lens of both correlations and regression coefficients. Average correlations over the full
sample provide evidence on hedge and diversifier roles, whereas rolling correlations and
tail–focused analyses help assess whether gold and silver behave as safe havens during
extreme market conditions. The linear models developed in the econometric sections further
quantify how strongly precious–metal returns respond to shocks in bond yields, equity
indices, and the US dollar, conditioning on multiple drivers at the same time. This
conceptual mapping between portfolio–theory definitions and econometric tools guides our
interpretation of the empirical results.
\subsection{Gold vs Silver}

Although gold and silver are often grouped together as ``precious metals'', their economic
roles and market dynamics exhibit important differences. Both metals share some safe–asset
characteristics, but their demand structures, industrial uses, and historical monetary
functions are not identical. Understanding these contrasts is crucial for interpreting the
empirical results on co–movements and for assessing whether findings for gold can be
generalized to silver.

From a historical perspective, gold has played a central role in international monetary
systems, serving as the ultimate anchor for currencies under the classical Gold Standard
and later as the reference asset in the Bretton Woods system. Silver, while also used as
money in various periods and regions, has progressively lost its monetary dominance and
has become more closely tied to industrial applications and jewelry demand. As a result,
gold tends to be perceived primarily as a monetary and financial asset, whereas silver is
more exposed to the business cycle and sector–specific industrial conditions.

In modern financial markets, gold is widely held by central banks, institutional
investors, and retail investors as a store of value and a potential hedge against
macroeconomic and financial risks. Silver, by contrast, has a smaller official–sector
presence and a larger share of demand coming from industrial uses, including electronics,
solar panels, and various manufacturing processes. This difference in demand composition
suggests that silver prices may be more sensitive to global growth expectations and
industrial activity, while gold may respond more strongly to monetary policy, real
interest rates, and currency developments.

These structural differences have implications for the hedge and safe–haven properties of
the two metals. Empirical studies often report that gold exhibits more robust hedge and
safe–haven behavior relative to equities and, in some cases, to currencies and bonds,
especially during periods of financial stress and heightened uncertainty
\sourcecitation{BaurLucey2010,ErbHarvey2013}. Silver, on the other hand, is sometimes found to offer
weaker or less stable protection, with its performance during crises depending more on the
specific nature of the shock and its impact on industrial demand.

At the same time, gold and silver are strongly related through common supply–side factors
and shared investor perceptions as precious metals. Their returns often display positive
co–movements over medium– to long–term horizons, reflecting joint responses to broad risk–
on and risk–off regimes, inflation concerns, and changes in commodity–market sentiment.
This combination of shared and distinct drivers makes the gold–silver relationship an
interesting test case for our econometric framework: if both metals were purely financial
safe assets, their behavior relative to bonds, equities, and the US dollar should be
similar; if industrial and cyclical forces matter more for silver, we should observe
systematic differences in correlations and regression coefficients.

In this paper we treat gold as the primary benchmark for safe–asset properties and use
silver as a comparative asset to disentangle monetary and industrial components in the
behaviour of precious–metal returns. The econometric analysis in later chapters will
quantify both the strength of the gold–silver co–movement and the extent to which their
relationships with bonds, equities, and the dollar differ across market regimes.
\subsection{Main Macroeconomic Drivers}

The behaviour of gold and silver cannot be understood solely by looking at their historical
monetary role or their classification as hedges and safe havens. Their prices are shaped by
a set of underlying macroeconomic and financial drivers that influence both investor demand
and opportunity costs. This subsection provides a conceptual overview of the main forces
that are most relevant for our analysis: real interest rates, inflation and inflation
expectations, the US dollar, equity–market risk, and global growth conditions.
\subsubsection{Real interest rates and opportunity cost}

Real interest rates play a central role in the valuation of non–yielding assets such as
gold. From a portfolio perspective, holding gold entails giving up the interest income that
could be earned on bonds or deposits. When real interest rates are high, the opportunity
cost of holding gold increases, which tends to exert downward pressure on its price. When
real rates are low or negative, the relative attractiveness of gold improves, as the
foregone yield is small or even negative in real terms.

This mechanism is particularly important in the post–Bretton Woods era, where gold is no
longer formally linked to currencies but remains a store of value whose demand is sensitive
to monetary policy and real yields. Silver, while also affected by interest–rate
conditions, may exhibit a weaker direct link to real rates because a larger share of its
demand is driven by industrial uses rather than pure store–of–value considerations.
\subsubsection{Inflation and inflation expectations}

Inflation and inflation expectations are another key driver of precious–metal demand. Gold
is often viewed as a hedge against unexpected inflation: when investors fear that the real
value of money and nominal assets will be eroded, they may shift into real assets such as
gold. This leads to a positive association between inflation scares and gold prices, even
if the relationship is not always stable across periods or regimes.

Silver can also benefit from inflation hedging demand, but its industrial component makes
its response more nuanced. In some episodes, higher inflation associated with strong
economic growth may support both industrial demand and silver prices; in others, inflation
shocks that weaken growth expectations can have mixed effects. Distinguishing between gold
and silver along this dimension is part of the motivation for treating silver as a
comparative asset in our empirical analysis.
\subsubsection{US dollar and exchange–rate conditions}

Because gold and silver are predominantly priced in US dollars on global markets, the
strength of the dollar is a fundamental determinant of their quoted prices. A stronger US
dollar tends to depress the prices of dollar–denominated commodities when expressed in
dollars, both through mechanical valuation effects and through shifts in international
demand. Conversely, a weaker dollar is often associated with higher precious–metal prices.

From the perspective of non–US investors, the dollar also acts as a competing safe asset.
In periods when the dollar appreciates strongly due to flight–to–quality flows or monetary
tightening, gold may face competition as a safe store of value. The interplay between
dollar strength and gold demand is therefore a crucial channel for understanding their
joint dynamics. Silver is influenced by similar currency effects, but again with a
potentially different sensitivity due to its industrial use and smaller role in official
reserves.
\subsubsection{Equity–market risk and risk sentiment}

Equity markets provide a natural proxy for global risk appetite and macroeconomic
conditions. In risk–on environments, characterized by rising stock prices and low risk
premia, investors may have less incentive to hold safe assets such as gold. In risk–off
episodes, marked by sharp equity sell–offs and elevated volatility, demand for safe stores
of value typically increases.

Within this framework, gold is often hypothesized to act as a hedge or safe haven against
equity risk, as discussed in the previous subsections. Silver, by contrast, has a more
ambiguous relationship with equities: its industrial component links it positively to
growth and corporate profits, while its precious–metal status may induce safe–asset
behaviour in severe crises. The net effect is an empirical question that our correlation
and regression analysis aims to address.
\subsubsection{Global growth and industrial demand}

Finally, global growth conditions and industrial activity play a more prominent role for
silver than for gold. Strong economic growth tends to increase industrial demand for
silver, supporting its price, while recessions and downturns can reduce usage in
manufacturing and technology applications. Gold, whose demand is more concentrated in
investment, jewelry, and official reserves, is less directly tied to industrial production
and may even benefit from growth slowdowns if they trigger risk aversion and accommodative
monetary policy.

This asymmetry implies that growth–related shocks can affect gold and silver in different
ways: a positive growth surprise may be neutral or slightly negative for gold (via higher
real rates and risk appetite) but positive for silver (via stronger industrial demand),
whereas a negative growth shock may boost gold as a safe asset while weighing on silver’s
industrial component. These patterns provide an additional dimension along which the
behaviour of the two metals can diverge.
\subsubsection{Implications for the empirical framework}

The macroeconomic drivers discussed above motivate the choice of explanatory variables in
our empirical framework. Real interest rates and bond yields capture the opportunity cost
channel; equity indices proxy for risk appetite and growth expectations; the US dollar
index summarizes currency and external–value effects; and growth–related variables
influence industrial demand, especially for silver. In subsequent chapters we translate
these conceptual drivers into specific time–series for bonds, equities, and the dollar,
and we investigate how gold and silver returns co–move with these factors across different
market regimes.

\section{The Post-COVID Repricing Context}

The contemporary relevance of gold cannot be reduced to the mechanics of the
classical gold standard. The post-COVID recovery combined supply constraints,
unusually expansionary policy settings and a rapid rebound in demand. The war
in Ukraine then intensified energy and food shocks while weakening global
growth prospects, reinforcing an inflationary environment already under way
\cite{imf-war-inflation-2022}. Subsequent geopolitical conflicts, sanctions,
trade fragmentation and uncertainty over the future path of public debt made
the monetary role of gold newly salient to investors and reserve managers.

These events do not identify a single stable law for the gold price. They alter
several channels at once. Inflation and inflation expectations affect the
perceived purchasing-power role of gold; real interest rates change the
opportunity cost of holding a non-yielding asset; the US dollar changes the
price faced by non-dollar investors; fiscal and monetary risk affect confidence
in nominal claims; reserve diversification affects official demand;
geopolitical risk changes precautionary demand; and mine supply, recycling,
jewellery and technology demand continue to determine the physical balance.
The empirical problem is therefore multivariate and regime dependent.

\section{The Debasement Trade: Narrative, Mechanisms and Limits}

The expression \emph{debasement trade} is used in market commentary for a
portfolio narrative in which investors seek scarce or non-sovereign assets
when they fear that inflation, fiscal dominance, currency depreciation or
monetary accommodation will erode the real value of nominal government money.
Gold is often placed at the centre of this narrative because it is not the
liability of a state or company and because it has a long monetary history.
That intuition is economically meaningful, but it is not sufficient evidence
that every rise in gold is caused by currency debasement.

Central-bank demand illustrates both the force and the ambiguity of the
narrative. Official gold accumulation can reflect reserve diversification,
liquidity, sanctions risk, geopolitical alignment or a desire to reduce
portfolio concentration; it need not imply an imminent rejection of the US
dollar. Federal Reserve research explicitly examines gold purchases together
with the continuing role of the dollar in international reserves, while ECB
analysis stresses the interaction between gold demand and geopolitical risk
\cite{fed-gold-reserves-2025,ecb-gold-risk-2025}. The World Gold Council's
current demand reporting provides a market-side account of official-sector
activity, but its classifications are descriptive rather than a causal test
\cite{wgc-gdt-q2-2026}.

For this thesis, the debasement trade is therefore a hypothesis-generating
framework. It motivates attention to regimes, tails and jumps, but it does not
replace econometric identification. Real rates, inflation, the dollar,
fiscal-monetary risk, reserve diversification, geopolitics and physical
demand/supply remain separate mechanisms. Their joint influence helps explain
why a valuation based only on a smooth extrapolation of past gold prices may
miss changes in market expectations and discontinuous repricing.

\section{From the Gold Narrative to the Barrick Valuation Object}

Starting Finance Research has a strong interest in corporate valuation. Gold
therefore becomes economically useful in this project not as an isolated trade
but as a stochastic revenue driver for a mining company. Barrick Mining
Corporation provides the case study through which commodity-price uncertainty
can be linked to production, costs, free cash flow and enterprise value. This
choice joins the historical and macroeconomic interpretation developed here to
the operating perimeter introduced in Chapter~3.

The next analytical step is deliberately postponed. Chapter~7 returns to the
legacy daily sample and asks what correlations, regressions, AR/ARMA/ARIMA/
ARFIMA diagnostics and conditional-volatility models can establish about gold
returns. Keeping that material out of the present chapter prevents historical
narrative from being confused with model fit.
